%! TEX root = science_techno_park.tex
% the main TeX file which is intended to compile, :VimtexReload after adjustment
   % this file should be included in the main file
% :h vimtex-tex-root
\documentclass[../wastek.tex]{subfiles}
%ensure the class of docs

\begin{document}

\section{Science Techno Park}%
\label{sec:Science Techno Park}

\subsection{Pendahuluan STP}%
\label{sub:Pendahuluan STP}


Perkembangan ilmu pengetahuan dan teknologi di bidang pertanian, pangan, dan biosains memiliki peran penting dalam meningkatkan ketahanan pangan serta daya saing suatu negara. Perguruan tinggi dan lembaga penelitian sebagai pusat inovasi sering menghasilkan berbagai temuan, namun belum seluruhnya dapat dimanfaatkan secara optimal oleh masyarakat dan industri.

Salah satu solusi untuk menjembatani kesenjangan antara penelitian dan penerapan adalah melalui pengembangan Science Techno Park (STP). Kawasan ini berfungsi sebagai pusat integrasi antara kegiatan riset, pengembangan teknologi, dan komersialisasi produk inovatif.

Pada dasarnya, persoalan yang dihadapi Indonesia bertumpu pada rendahnya hasil riset dan teknologi yang dapat diadopsi oleh industri. Hal ini disebabkan oleh ketidakpaduan antara kebutuhan industri dengan teknologi yang dikembangkan, sehingga banyak inovasi yang tidak dapat diimplementasikan oleh masyarakat (Sujarwo, 2012).

Untuk menjembatani permasalahan tersebut, pemerintah melalui Nawa Cita ke-6 berupaya membangun sejumlah science and techno park (STP) skala nasional, serta mengembangkan politeknik dan SMK dengan prasarana, sarana, dan teknologi terkini (International Association of Science Parks, 2015).

Pengembangan STP di Indonesia diatur dalam Undang-Undang Nomor 18 Tahun 2002 tentang Sistem Nasional Penelitian, Pengembangan, dan Penerapan Ilmu Pengetahuan dan Teknologi. Selanjutnya, kebijakan ini diperkuat melalui Undang-Undang Nomor 106 Tahun 2017 tentang Kawasan Sains dan Teknologi (KST). Regulasi tersebut mendorong berbagai pihak, khususnya pemerintah, untuk membangun kawasan yang berfungsi sebagai sarana dan prasarana pengembangan serta transfer ilmu pengetahuan dan teknologi (Wibowo, 2017).

Pembangunan STP bertujuan untuk memadukan hasil riset dan teknologi agar dapat diadopsi oleh industri serta dimanfaatkan oleh masyarakat. Selain itu, penguatan STP juga dilatarbelakangi oleh skor Global Innovation Index (GII) Indonesia pada tahun 2021 sebesar 27,1 (peringkat 87 dari 126 negara), yang mengalami penurunan dibandingkan dua tahun sebelumnya (Rosa, 2022).

Secara konseptual, STP merupakan kawasan yang diperuntukkan bagi penelitian dan pengembangan sains dan teknologi, yang dikelola secara profesional melalui penciptaan dan penguatan ekosistem inovasi. Tujuannya adalah untuk meningkatkan daya saing industri atau institusi yang berada di dalamnya, sekaligus mendorong pertumbuhan ekonomi yang berkelanjutan. Ciri utama STP adalah adanya proses transfer teknologi dan komersialisasi produk inovasi (Sari \& Retnaningsih, 2020).

Lebih lanjut, STP berperan dalam mengelola arus pengetahuan dan teknologi antara universitas, lembaga penelitian dan pengembangan, serta industri. STP juga memfasilitasi penciptaan dan pertumbuhan perusahaan berbasis inovasi melalui inkubasi bisnis, spin-off, serta penyediaan ruang dan fasilitas pendukung lainnya (Soenarso et al., 2013).

Pembangunan STP di Indonesia yang cukup masif, termasuk yang dikelola oleh pemerintah daerah, tidak boleh hanya menjadi tren semata. Diperlukan penguatan kualitas, khususnya pada STP yang telah terbentuk di tingkat daerah. Hingga saat ini, terdapat sekitar 10 STP yang dikelola oleh pemerintah daerah. Di antaranya, lima STP yang menjadi objek perbandingan adalah Cimahi Techno Park, Solo Techno Park, Sragen Techno Park, STP Sumatera Selatan, dan Riau STP.

Secara umum, STP tersebut masih dalam tahap pencarian bentuk terkait pola pengelolaan dan pengembangannya. Oleh karena itu, diperlukan perumusan pola kelembagaan STP pemerintah daerah yang tepat, guna mempercepat inovasi teknologi melalui sinergi antara lembaga pendidikan, industri, dan komunitas. Selain itu, STP juga berpotensi menciptakan lapangan kerja baru serta meningkatkan pendapatan negara melalui pajak (Seo, 2006).

Dengan demikian, diperlukan strategi pengembangan STP yang komprehensif, khususnya bagi STP yang dikelola oleh pemerintah daerah, agar dapat berfungsi secara optimal dalam mendorong inovasi dan pertumbuhan ekonomi.


\subsection{Studi Kasus Science Techno Park di Indonesia}%
\label{sub:studi_kasus_science_techno_park_di_indonesia}
% Tambahkan di preamble

\tabitem Cimahi Techno Park
Cimahi Techno Park merupakan Unit Pelaksana Teknis (UPT) di bawah Dinas Perdagangan, Koperasi, Usaha Kecil dan Menengah, serta Perindustrian Kota Cimahi. STP ini bertujuan sebagai wadah pengembangan dan pemanfaatan ilmu pengetahuan dan teknologi untuk mendorong pertumbuhan ekonomi kreatif.

Cimahi Techno Park mengembangkan 16 subsektor ekonomi kreatif, yaitu:
\begin{enumerate}
	\item aplikasi dan game developer
	\item arsitektur
	\item desain interior
	\item desain komunikasi visual
	\item desain produk
	\item fashion
	\item film, animasi dan video
	\item fotografi
	\item kriya
	\item kuliner
	\item musik
	\item penerbitan
	\item periklanan
	\item seni pertunjukan
	\item seni rupa
	\item televisi dan radio
\end{enumerate}

Fokus utama pengembangannya meliputi kolaborasi riset, inkubasi bisnis dan spin-off, serta layanan teknologi.

\tabitem Solo Techno Park
Solo Techno Park pada awalnya merupakan UPT di bawah Bappeda Kota Surakarta, kemudian pada tahun 2010 berubah menjadi Badan Layanan Umum Daerah (BLUD) untuk meningkatkan fleksibilitas pengelolaan.

Tantangan utama yang dihadapi meliputi sulitnya menjalin kerja sama dengan industri, penentuan tarif layanan, serta perbedaan pola pikir antara pemerintah dan dunia industri.

Sebagai solusi, Solo Techno Park difokuskan sebagai penyedia tenaga kerja terampil yang siap digunakan oleh industri, khususnya industri mitra.

\tabitem Sragen Techno Park
Sragen Techno Park diresmikan pada tahun 2009 dengan fokus pada pengembangan tenaga kerja berbasis kompetensi.

Program utama yang dijalankan meliputi pelatihan berbasis Standar Kompetensi Kerja Nasional Indonesia (SKKNI) dan inkubasi bisnis berbasis teknologi.

Bidang pelatihan mencakup:
\begin{enumerate}
	\item otomotif
	\item pengelasan
	\item industri tekstil
	\item konstruksi bangunan
	\item listrik
	\item mekanik logam
	\item tata niaga
\end{enumerate}

Konsep utama yang diterapkan adalah One Stop Service Labor Market, yaitu penyediaan tenaga kerja siap pakai bagi industri.

\tabitem STP Sumatera Selatan
STP Sumatera Selatan merupakan pengembangan dari Agroteknopark Palembang yang berfokus pada integrasi sektor pertanian, perikanan, peternakan, dan pascapanen secara terpadu.

Fungsi utama kawasan ini meliputi pusat alih teknologi, pusat percontohan pertanian berbasis teknologi, serta pengembangan inovasi pertanian.

Inovasi yang dikembangkan antara lain:
\begin{enumerate}
	\item rekayasa genetik ternak
	\item teknologi pakan berbasis sumber daya lokal
	\item rekayasa iklim mikro
	\item sistem rantai pasok dingin
\end{enumerate}

Tantangan yang dihadapi meliputi keterbatasan investasi dan rendahnya tingkat adopsi teknologi oleh industri.

\tabitem Riau Science Techno Park
Riau Science Techno Park berada di bawah Bappeda dan dikembangkan di tiga wilayah, yaitu Kabupaten Kampar, Pelalawan, dan Siak.

Fokus pengembangan di masing-masing wilayah:
\begin{enumerate}
	\item Kampar: pangan lokal (sagu, kelapa, nenas, ikan)
	\item Pelalawan: industri hilir kelapa sawit
	\item Siak: benih padi, hortikultura, dan peternakan
\end{enumerate}

Produk unggulan yang dihasilkan antara lain:
\begin{enumerate}
	\item pindang patin kaleng
	\item kerupuk ikan
	\item selai nenas
\end{enumerate}

Kendala yang dihadapi meliputi lokasi yang sulit dijangkau, keterbatasan sumber daya manusia, serta hambatan dalam transfer teknologi.

\subsection{Strategi Pengembangan Science Techno Park}%
\label{sub:strategi_pengembangan_science_techno_park}


Berdasarkan studi kasus, terdapat beberapa komponen utama dalam pengembangan STP, yaitu:

\begin{enumerate}
	\item infrastruktur dan lokasi yang strategis
	\item manajemen yang profesional
	\item sumber pengetahuan dari perguruan tinggi atau lembaga riset
	\item tenant atau startup berbasis inovasi
	\item inkubator bisnis
	\item keterlibatan industri
\end{enumerate}

Namun, dalam praktiknya banyak STP masih terbatas pada pelatihan tenaga kerja dan pembinaan UMKM, sehingga belum sepenuhnya memenuhi konsep STP yang ideal.

\tabitem Model Kelembagaan BLUD
BLUD merupakan model pengelolaan yang memungkinkan kolaborasi antara pemerintah, swasta, dan akademisi. Model ini lebih fleksibel dibandingkan UPT.

Keunggulan BLUD meliputi fleksibilitas pengelolaan, efisiensi, serta kemampuan menghimpun pembiayaan dari berbagai pihak.

Namun, tantangan yang dihadapi adalah perbedaan kepentingan antar aktor dan keterbatasan dana operasional.

\tabitem Model Kolaborasi Pentahelix
Pengembangan STP membutuhkan kolaborasi lima aktor utama, yaitu academia, business, government, community, dan media.

\tabitem Prinsip Kolaborasi Efektif
Agar pengembangan STP berjalan optimal, diperlukan:
\begin{enumerate}
	\item pembagian tugas yang jelas
	\item kesepakatan antar aktor
	\item mekanisme koordinasi yang baik
	\item sistem penyelesaian konflik
	\item manajemen profesional
	\item keseimbangan antara kepentingan publik dan bisnis
\end{enumerate}




\bibliographystyle{apalike}
\bibliography{../filename.bib} % required for citecompletion, not work in mainfile
\end{document}

